\section{Method}
\label{sec:method}
\paragraph{Verification.} A frozen CNN $f$ turns an image into an embedding of length one. For an anchor image $a$ and a probe image $p$, the score is the cosine similarity $s(a,p)=f(a)^\top f(p)$. If $s\ge\tau$ we say ``same person''. We pick the threshold $\tau$ on validation pairs only, choosing the value with the highest accuracy (lowest error). As a check, we also fit a Gaussian to the genuine scores and another to the impostor scores and compute the Bayes threshold. Both methods agree within about 0.05 ($\tau\approx0.42$--$0.48$ across defenses).

\paragraph{Threat model and attacks.} The attacker can only change the probe: $p'=p+\delta$ with $\|\delta\|_\infty\le\epsilon$ and pixel values kept in $[0,1]$. The result is then saved as a normal 8-bit image. For a genuine pair (label $y{=}1$) the attacker wants the model to say ``different'' (\emph{dodging}, lower the score). For an impostor pair ($y{=}0$) the attacker wants ``same'' (\emph{impersonation}, raise the score). Using $\sigma=+1$ for impostors and $-1$ for genuine pairs, both goals become one problem:
\begin{equation}
\max_{\|\delta\|_\infty\le\epsilon}\; \sigma\, s(a,p+\delta).
\end{equation}
PGD repeats $p_{k+1}=\Pi\big(p_k+\alpha\,\mathrm{sign}(\nabla_{p}\,\sigma s)\big)$ from a random starting point, and FGSM is the same thing done in one full step. DeepFool takes the smallest $\ell_2$ step that crosses $\tau$, and CW minimizes $\|\delta\|_2^2+c\cdot\max(\cdot)$ in tanh space; we clip both to the $\ell_\infty$ limit. Square Attack only uses the scores. These attacks are first run \emph{non-adaptively}, i.e.\ they are made against the model with no defense.

\paragraph{Adaptive attacks.} Many filters $t$ have no useful gradient. BPDA gets around this by using $t(p)$ in the forward pass and pretending it is the identity in the backward pass: $\tilde t(p)=p+\mathrm{sg}(t(p)-p)$, where $\mathrm{sg}$ means ``stop gradient''. Against a single filter, the attacker maximizes $\sigma s(t(a),t(p'))$. Against our system, the \emph{ensemble} attack takes the average of this over no filter plus all five of our filters. The \emph{detector-aware} version goes one step further and also maximizes $\cos\big(f(p'),f(t(p'))\big)$ for every $t$. In plain words, it tries to keep the embedding steady under the filters, which is exactly the signal our detector looks for.

\paragraph{Consistency features.} We use five filters $\mathcal T=\{$JPEG-75, JPEG-30, median $3{\times}3$, Gaussian, 4-bit depth$\}$. Let $s_0=s(a,p)$ be the normal score and $s_t=s(t(a),t(p))$ the score after filter $t$. The feature vector is
\begin{equation}
\begin{split}
z&=\Big[s_0,\ \big(s_t-s_0,\ 1-\min\{c^t_a,\,c^t_p\}\big)_{t\in\mathcal T}\Big],\\
c^t_x&=\cos\big(f(x),f(t(x))\big),
\end{split}
\end{equation}
which has $1{+}2|\mathcal T|=11$ numbers. The idea is simple. Attack noise is tuned very finely to the raw model, so it is fragile. When we filter the image, the noise gets damaged, the score jumps and the embedding moves (``drifts''). A clean pair, on the other hand, barely changes.

\paragraph{Detector and Bayes fusion.} We train four standard classifiers on $z$ to tell clean pairs from attacked ones: LDA (Gaussians with a shared covariance), Gaussian Bayes with a separate covariance per class (QDA), an SVM with an RBF kernel, and boosted shallow trees. All use equal class priors. The detector gives a probability $P(\mathrm{adv}\mid z)$. We combine it with the score using the minimum-risk rule: flag the pair if $P(\mathrm{adv}\mid z)\,C_{\text{miss}}>P(\mathrm{clean}\mid z)\,C_{\text{fa}}$, which is the same as $P(\mathrm{adv}\mid z)>C_{\text{fa}}/(C_{\text{fa}}+C_{\text{miss}})$. If the pair is not flagged, we answer Same or Different using $s_0\ge\tau$. In the experiments we use a threshold that is easier to understand: we set it so that 5\% of clean \emph{validation} pairs get flagged. This gives $0.303$ for the chosen model (for comparison, the cost-based value with $C_{\text{miss}}{:}C_{\text{fa}}=5{:}1$ would be $0.167$).
